\documentclass[letterpaper]{article}
\usepackage[submission]{aaai2027}
\usepackage[hyphens]{url}
\usepackage{graphicx}
\urlstyle{rm}
\def\UrlFont{\rm}
\usepackage{natbib}
\usepackage{caption}
\frenchspacing
\usepackage{amsmath}
\usepackage{amssymb}
\usepackage{booktabs}
\usepackage{algorithm}
\usepackage{algorithmic}

\pdfinfo{
/TemplateVersion (2027.1)
}

\setcounter{secnumdepth}{1}

% Working title. Do not freeze until the inverse-dynamics comparison and
% task-relevant effect metric are complete.
\title{Making Learned Dynamics Actionable: Posterior Control-Equivalent
Transport for Source-Only Continual Adaptation}

\author{Anonymous Submission}
\affiliations{}

\begin{document}

\maketitle

\begin{abstract}
Policies trained under one set of physical dynamics can fail when mass,
gravity, friction, or actuation changes at deployment.  An online world model
can relearn the changed transition law, but accurate prediction alone does not
specify how a frozen task policy should change its action.  We study this
\emph{actionability gap} in a source-only continual adaptation setting: the
policy and dynamics model are trained in one source environment, target
dynamics are unlabeled and may recur, and deployment provides only observed
state--action transitions.  We introduce posterior control-equivalent
transport (PCET), which interprets a source-policy action through its predicted
local physical effect and pulls that effect back to an action under the current
online dynamics posterior.  A Bayesian KAN model supplies the local drift,
action-effect operator, and directional uncertainty required by a closed-form,
risk-regularized pullback; no target-domain policy optimization or online
trajectory search is required.  A non-overwriting memory indexed by
control-equivalent residuals supports recurring dynamics.  Across a two-link
nonlinear system and continuous-action CartPole, the current implementation
recovers frozen policies from severe dynamics shifts using only self-supervised
transition updates.  The results indicate that the primary benefit is faster
performance recovery rather than merely a higher asymptotic success rate.
\end{abstract}

\section{Introduction}

Reinforcement-learning policies often entangle a task---such as balancing or
reaching---with the particular dynamics under which the policy was trained.
Consequently, a policy that succeeds in a nominal environment may fail after a
change in payload, gravity, link inertia, friction, or actuator strength.  A
common response is to expose the policy to a distribution of dynamics during
training or to adapt the policy itself after deployment.  The former requires
anticipating a useful variation family, while the latter consumes task reward
and can damage previously learned behavior.

World models offer a different source of adaptation.  Recent transitions can
identify how the current system responds to actions, without a physical
parameter label or a task-success signal.  Context-aware dynamics models
\cite{lee2020cadm}, augmented world models \cite{ball2021augwm}, and online
world-model planning \cite{liu2025onlinewm} demonstrate the value of adapting
the predictive model.  However, an accurate estimate of
$F_t(s,a)$ answers ``what will this action do?'' rather than ``which action now
preserves the task intended by a source policy?''  We call this missing map
from updated dynamics knowledge to action correction the \emph{actionability
gap of learned dynamics}.

Matching action effects across domains is not itself new.  Deep inverse
dynamics transfer predicts the state change intended by a source controller
and inverts a target model to reproduce it \cite{christiano2016inverse};
Grounded Action Transformation (GAT) similarly learns action transformations
for simulator grounding \cite{hanna2021gat}.  Task--dynamics decoupling has
also been used to convert high-level task commands to platform-specific control
signals \cite{liu2023dttd}.  Our question is narrower and continual: can a
frozen task policy be re-grounded from a \emph{source-only}, self-supervised
forward-dynamics posterior when unannounced dynamics changes recur, without a
separately trained black-box inverse model, target reward optimization, or
online model-predictive search?

Our answer is posterior control-equivalent transport (PCET).  A source action
$a^{\mathrm{src}}=\pi_0(s)$ is first encoded as the local physical effect that
the source model predicts it should cause.  A continually updated Bayesian KAN
represents the current dynamics as a local drift plus an action-effect operator.
PCET then solves a risk-regularized inverse problem that reproduces the source
effect under the current dynamics.  Under a locally control-affine model this
inverse has a closed form, and every executed action depends explicitly on the
current drift, action-effect matrix, and posterior risk.  The world model is
therefore on the mandatory decision path rather than an optional context that
the policy can ignore.

We make three contributions.  First, we formalize source-only continual
dynamics adaptation as an actionability problem, separating task knowledge in
a frozen policy from its changing physical realization.  Second, we derive a
posterior-risk control-equivalent pullback from an online forward model to the
action space, together with a structured KAN realization exposing the local
action-effect operator.  Third, we combine the pullback with
control-relevant posterior stabilization and non-overwriting dynamics memory,
and evaluate recovery and retention under multiple unseen dynamics shifts.

\section{Related Work}

\paragraph{Dynamics generalization and context.}
CaDM infers a latent dynamics context from recent experience and conditions a
forward model on it \cite{lee2020cadm}.  AugWM starts from a single offline
environment but constructs an explicit augmentation family during training and
conditions the policy on the inferred augmentation \cite{ball2021augwm}.
These methods show that recent transitions contain useful dynamics information.
PCET instead asks how an online forward posterior can alter the physical action
of an already trained policy through an explicit operator.

\paragraph{Action grounding and task--dynamics separation.}
Inverse-dynamics transfer and GAT provide the closest conceptual precedents to
effect matching \cite{christiano2016inverse,hanna2021gat}.  DTTD separates a
task-oriented reinforcement-learning module from a low-level dynamics
controller in a UAV interception system \cite{liu2023dttd}.  A recent dual
control formulation also treats world-model predictions as reference
trajectories for rapid motor adaptation \cite{brito2025rwm}.  PCET does not
claim task--dynamics separation or effect matching as new in isolation.  Its
candidate contribution is the posterior, forward-model-based and continual
formulation: uncertainty affects the inverse operator directionally, only the
cognitive model is updated online, and recurring dynamics are recalled without
overwriting other posterior states.

\paragraph{Continual world models.}
Planning with online world models can provide continual decision making and
formal regret guarantees under suitable assumptions \cite{liu2025onlinewm}.
Our setting retains a learned task policy and avoids online action-sequence
planning.  It focuses on the interface that re-realizes that policy under a
changing plant.  KANs \cite{liu2024kan} are used here as a structured local
function dictionary, not on the assumption that spline networks are inherently
immune to catastrophic forgetting.

\section{Problem Formulation}

Let $\mathcal{M}_{\xi}=(\mathcal{S},\mathcal{A},P_{\xi},r,\gamma)$ denote a
family of continuous-control MDPs.  Across the family, state semantics, action
semantics, reward, and the task remain fixed; only the transition dynamics
$P_{\xi}$ change.  Training data are available from one source dynamics
$\xi_0$ only.  Source training produces a task policy $\pi_0$ and a source
dynamics model $F_0$.  Both may use any valid source-domain optimization, but
their losses are separate.

At deployment, the unobserved dynamics index $\xi_t$ may change abruptly or
recur.  The agent receives neither $\xi_t$ nor a change-point notification.  It
observes only transition tuples
\begin{equation}
\mathcal{D}_t=\{(s_i,a_i,s_{i+1})\}_{i=1}^{N_t}.
\end{equation}
The policy parameters remain frozen.  Only the dynamics posterior and its
memory may update, using a one-step prediction objective.  No target reward,
teacher action, or target-domain policy gradient is used.

We distinguish three evaluation regimes.  \emph{Zero-shot transfer} is the
performance immediately after a shift and before any target update.
\emph{Recovery} is performance as a function of the number of target
transitions.  \emph{Retention} is the performance and recovery delay when a
previously encountered dynamics condition returns.  The main objective is the
area under the recovery curve, subject to a small final performance gap and
limited forgetting, rather than an unidentifiable guarantee of zero-shot
success for arbitrary unseen physics.

\section{Posterior Control-Equivalent Transport}

\subsection{Local Physical Effects}

Let $y\in\mathbb{R}^{d}$ be a local physical effect derived from a transition,
such as acceleration or a selected state increment.  The cognitive model is
\begin{equation}
y=F_{\theta}(s,a)+\epsilon.
\end{equation}
We use a locally control-affine representation
\begin{equation}
F_{\theta}(s,a)=f_{\theta}(s)+G_{\theta}(s)a,
\quad G_{\theta}(s)\in\mathbb{R}^{d\times m}.
\label{eq:affine}
\end{equation}
This is an exact model for control-affine systems and a local Jacobian model for
general differentiable action dynamics.  It does not require naming mass,
gravity, or friction; $f$ and $G$ are identifiable function-level quantities
needed for control.

For the frozen source action
\begin{equation}
a^{\mathrm{src}}=\pi_0(s),
\end{equation}
the source model defines the intended effect
\begin{equation}
v_0(s)=F_0(s,a^{\mathrm{src}}).
\label{eq:intent}
\end{equation}
Equation~\eqref{eq:intent} does not assert that every source next state must be
copied.  More generally, an effect metric $W(s)\succeq 0$ can retain only
task-relevant directions.

\subsection{Effect Geometry and Automaticity Boundary}

The current reported method uses $W=I$.  A task-dependent $W(s)$ can improve a
chosen benchmark when its directions are specified manually, but such a result
does not establish automatic transfer.  We therefore exclude manually weighted
metrics from the empirical claims.  A source-only value Critic predicts returns
accurately in our diagnostic study, yet both its gradient outer product and its
positive-semidefinite value Hessian fail as effect metrics.  This indicates that
a local symmetric metric is insufficient to represent the full, asymmetric
counterfactual advantage of a physical effect.  Learning and transporting that
advantage is left as an open extension rather than hidden behind a hand-crafted
weight.

\subsection{Bayesian KAN Dynamics Posterior}

Let $\phi(s)$ be a fixed multiresolution KAN dictionary containing local
B-spline features and low-order cross-state interactions.  Action-conditioned
features are
\begin{equation}
\Phi(s,a)=\left[\phi(s),\frac{a_1}{\sigma_1}\phi(s),\ldots,
\frac{a_m}{\sigma_m}\phi(s)\right],
\end{equation}
and the effect prediction is $\hat y=\Phi(s,a)B$.  Fixing the dictionary keeps a
stable coordinate system while its coefficients are updated recursively.
For a local chart $k$, sufficient statistics follow
\begin{align}
\Lambda_{k,t}&=\Lambda_{k,t-1}+X_{k,t}^{\top}X_{k,t},\\
b_{k,t}&=b_{k,t-1}+X_{k,t}^{\top}y_t,\\
\bar B_{k,t}&=\Lambda_{k,t}^{-1}b_{k,t},\quad
\Sigma_{k,t}=\Lambda_{k,t}^{-1}.
\end{align}
Local charts are blended by fixed-support weights and may grow when deployment
data reach an uncovered state region.  The resulting posterior supplies
$\bar f_t(s)$, $\bar G_t(s)$, and action-directional epistemic uncertainty.

Before updating with a new micro-batch, we compute the prequential residual
\begin{equation}
r_t=y_t-\hat y_{t\mid t-1}.
\end{equation}
Its control-equivalent form is
\begin{equation}
e_t^a=(\bar G_t^{\top}\bar G_t+\eta I)^{-1}\bar G_t^{\top}r_t,
\label{eq:action-residual}
\end{equation}
which measures how much action error would have an effect comparable to the
model residual.  Unlike raw prediction RMSE, this coordinate discounts errors
that the action cannot correct and preserves errors that matter strongly to
control.

\subsection{Posterior-Risk Pullback}

Let $U_t(s)\succ 0$ denote the action-space risk obtained by propagating the
KAN posterior covariance through the action-effect operator.  PCET selects the
current physical action by
\begin{equation}
\begin{aligned}
a_t^*=\arg\min_a\quad
&\left\|\bar f_t(s)+\bar G_t(s)a-v_0(s)\right\|_{W(s)}^2\\
&+\left\|a-a^{\mathrm{src}}\right\|_{U_t(s)}^2.
\end{aligned}
\label{eq:pullback-objective}
\end{equation}
For the local model in Equation~\eqref{eq:affine}, the solution is
\begin{equation}
\begin{aligned}
H_t={}&\bar G_t^{\top}W\bar G_t+U_t,\\
a_t^*={}&H_t^{-1}
\left[\bar G_t^{\top}W(v_0-\bar f_t)+U_ta^{\mathrm{src}}\right].
\end{aligned}
\label{eq:pullback}
\end{equation}
The first term reproduces the source effect under the new dynamics.  The second
shrinks toward the verified source action along uncertain or weakly
identifiable directions.  Equation~\eqref{eq:pullback} supports $d\neq m$ and
handles both under- and over-actuated local maps through regularization.

If the current model equals the source model, the source action is a minimizer.
More importantly, the executed action explicitly depends on
$(\bar f_t,\bar G_t,U_t)$ at every step, which rules out a policy-side shortcut
that ignores the cognitive model.  PCET performs one matrix solve rather than
rolling out a model or searching an action sequence.

\subsection{Stable Transfer and Recurring Dynamics}

Fast posterior updates can produce abrupt control changes before the relevant
operator directions are sufficiently identified.  For difficult shifts, we
maintain a fast cognitive posterior and a decision-side posterior.  Candidate
coefficient changes are decomposed under a control metric and admitted
continuously in the directions supported by evidence.  This spectral transport
is an online stabilization layer, not a prerequisite for the core pullback.

A single posterior still forgets when the same state--action region supports
conflicting transition functions.  We therefore maintain a set of
non-overwriting posterior contexts.  For each context $k$, recent transitions
produce control-equivalent residuals from Equation~\eqref{eq:action-residual}
and a Mahalanobis score
\begin{equation}
m_k=\frac{1}{N}\sum_i
\tilde e_i^{(k)\top}(C_k+\varepsilon I)^{-1}\tilde e_i^{(k)}.
\end{equation}
A sticky transition prior combines this score with temporal persistence.  An
existing posterior is recalled when it explains the batch; otherwise a new
posterior is created.  Updating one context does not overwrite the sufficient
statistics of another.

\begin{algorithm}[t]
\caption{Online Posterior Control-Equivalent Transport}
\label{alg:pcet}
\begin{algorithmic}[1]
\REQUIRE Frozen source policy $\pi_0$, source model $F_0$, posterior memory
\FOR{each control step $t$}
  \STATE Observe state $s_t$ and compute $a^{\mathrm{src}}=\pi_0(s_t)$
  \STATE Encode source intent $v_0=F_0(s_t,a^{\mathrm{src}})$
  \STATE Query current posterior for $\bar f_t,\bar G_t,U_t$
  \STATE Compute $a_t$ with Equation~\eqref{eq:pullback} and execute it
  \STATE Observe $s_{t+1}$ and form the local effect $y_t$
  \STATE Score, recall, or create a posterior context using $e_t^a$
  \STATE Update only the selected cognitive posterior with $(s_t,a_t,s_{t+1})$
\ENDFOR
\end{algorithmic}
\end{algorithm}

\section{Experiments}

We evaluate three questions: (i) whether an accurate target dynamics model can
recover a frozen source policy, (ii) how full-effect PCET compares with learned
inverse transfer under a shared data budget, and (iii) whether recurring dynamics can be recalled without catastrophic
performance loss.  The policy is trained under one source
dynamics only.  At deployment, no method receives the target physical
parameters or an environment identifier.

\subsection{Continuous-Action CartPole}

The CartPole state is $[x,\theta,\dot x,\dot\theta]$ and the action is a scalar
continuous force.  The source actor achieves 100\% success over 1,024 random
initial states, and the source KAN has acceleration RMSE $0.00488$.  We test
three shifts that reduce the frozen actor to 0\% success.  Table~\ref{tab:cartpole}
reports five-seed means after 2,048 target transitions.  All online methods see
identical target transition prefixes and test states; the learned inverse KAN
has 271 coefficients versus 300 for the forward KAN.

\begin{table}[t]
\centering
\small
\caption{CartPole final success under unseen dynamics (\%).  The actor remains frozen.}
\label{tab:cartpole}
\begin{tabular}{@{}lrr@{}}
\toprule
Shift & PCET & Inverse \\
\midrule
Heavy cart, weak act. & 99.96 & \textbf{100.00} \\
High gravity, weak act. & 98.79 & \textbf{99.08} \\
Hard combined shift & 78.63 & \textbf{79.51} \\
\bottomrule
\end{tabular}
\end{table}

For the hard shift, PCET and direct inverse transfer obtain mean recovery rates
of 71.45\% and 70.22\%, respectively, and neither dominates at the final
budget.  This strong comparison prevents us from attributing the present result
to a uniquely effective forward-model transport mechanism.  A manually
weighted diagnostic performs better, but is excluded because its task geometry
was not recovered automatically.

\subsection{Two-Link Nonlinear Control}

The two-link environment has four-dimensional state, two-dimensional action,
and two-dimensional acceleration effects.  We evaluate heavy-load, low-gravity,
and coupled multi-parameter shifts at budgets of 128 to 2,048 transitions over
five seeds.  Table~\ref{tab:twolink} compares the previous spectral-transport
baseline with posterior-conditioned decision making.

\begin{table}[t]
\centering
\small
\caption{Five-seed two-link recovery results (success rate, \%).}
\label{tab:twolink}
\begin{tabular}{@{}lrr@{}}
\toprule
Metric & Spectral & PCET \\
\midrule
Mean recovery & 75.89 & \textbf{82.76} \\
Final mean & 95.55 & \textbf{95.70} \\
Worst final & 71.48 & \textbf{72.66} \\
Coupled recovery & 44.69 & \textbf{60.31} \\
Coupled seeds at 95\% & 1/5 & \textbf{3/5} \\
\bottomrule
\end{tabular}
\end{table}

The largest gain is in the recovery trajectory rather than the final mean.  In
a separate three-seed sequence---heavy load, low gravity, coupled shift, and
then revisits---the first 128-transition revisit batch averages 98.87\%
success, the final revisit mean is 98.96\%, and the worst final revisit is
90.63\%.  The memory creates two duplicate contexts, so these results support
retention of control performance but not optimal memory allocation.

\subsection{Required Final Comparisons}

The current evidence establishes feasibility, cross-system recovery, and a
unified-budget learned-inverse comparison on CartPole.  The final empirical
claim still requires (i) an automatic source-value transfer operator, (ii) a context-conditioned actor, (iii) online policy
fine-tuning, (iv) model-predictive control, and (v) an equal-capacity MLP
forward model.  The task metric must also be validated on the two-link system,
and posterior risk, stable transport, and memory require individual ablations.

\section{Limitations}

The current experiments use simulated systems and safe exploratory transitions
with randomized resets.  CartPole does not yet include recurrent-environment
memory, and the two-link memory can create duplicate contexts.  Manually
weighted effect geometry is not counted as evidence; automatic Critic-gradient
and Critic-Hessian metrics currently underperform $W=I$.  Finally, local control affinity and sufficient excitation
are substantive assumptions: if the action-effect operator is poorly
conditioned or the target effect is infeasible, regularization can limit
unsafe corrections but cannot guarantee task recovery.

\section{Conclusion}

We study how an online dynamics model can become operationally useful to a
frozen task policy under unseen and recurring physical changes.  PCET encodes
the source policy through its intended local effect and uses a Bayesian KAN
posterior to compute an uncertainty-aware action that realizes that effect
under the current dynamics.  Results show source-only recovery in systems with
different action dimensions, but a unified-budget CartPole comparison finds no
clear advantage over learned inverse transfer.  The central remaining problem
is to transport a source task's full counterfactual value information without
manual effect weights or target-domain supervision.

\bibliography{references}

\end{document}
